\scanpage{10}
% Pre-ruled music paper.
dire que $L_{\bullet}\simeq L_{\bullet}^{\circ}\otimes_{\mathbb Z}k$, où $L_{\bullet}^{\circ}$ est un complexe de $\mathbb Z$-modules $\frac12$ simplicial spécial.

\uline{Proposition 2} Les \deleted{\unsure{$\mathbb Z$}}$k$-modules \inserted{$\frac12$ simpliciaux} spéciaux (resp. str. spéciaux) sont stables par produits tensoriels [et par les op. tensorielles $\Lambda^i,\Gamma^i,\operatorname{Sym}^i$ ?]
% slanted
\margin{\unsure{lemme}}
% slanted
\margin{[ ] faux}

Le cas strictm spécial se ramène au cas spécial, en se plaçant sur $\mathbb Z$\ldots

1°) Produit tensoriel. On est réduit à prouver que les modules
\[
\left\{\begin{aligned}
\Lambda^i\Phi_*^{\deleted{\unsure{*}}}\otimes\Lambda^j\Phi_* &\qquad(i,j\geq1)\\
\Lambda^i\Psi_*\otimes\Lambda^j\Psi_* &\qquad(i,j\geq0)\\
\Lambda^i\Phi_*\otimes\Lambda^j\Psi_* &\qquad(i\deleted{,j}\geq1,j\geq0)
\end{aligned}\right.
\]
sont spéciaux. Or ils sont à composantes libres de t.f., et par Künneth-Eilenberg-Zilber leurs $\pi_i$ sont itou, il reste à montrer qu'il n'y a qu'un nb fini de facteurs $M_i,N_j\ne0$. \deleted{On va} \inserted{Cela résulte} \deleted{prouver} aussitôt du

\deleted{\uline{Lemme} \deleted{Sous} Les conditions du cor. 2 de prop 1 sont encore équivalentes à \unsure{celles|les|la} suivantes (en supposant \inserted{$k$ comm. et} $\operatorname{Spec}k$ connexe, de sorte que $\mu_i=\operatorname{rg}M_i$ et $\nu_j=\operatorname{rg}N_j$ est défini)

a) $\operatorname{rg}_k L_n$ est un polynôme en $n$ pour $n$ grand

b) $\exists$}
% archive sheet number: 9
